\documentclass[11pt]{article}
\usepackage[a4paper,margin=22mm]{geometry}
\usepackage[T1]{fontenc}
\usepackage[utf8]{inputenc}
\usepackage{lmodern,microtype}
\usepackage{amsmath,amssymb,bm,booktabs,array,longtable,graphicx}
\usepackage{natbib,caption,subcaption,xcolor,hyperref,placeins}
\hypersetup{colorlinks=true,linkcolor=blue!55!black,citecolor=blue!55!black,urlcolor=blue!55!black}
\setlength{\emergencystretch}{2em}
\setlength{\parskip}{0.25em}
\setlength{\parindent}{1em}
\renewcommand{\topfraction}{.92}
\renewcommand{\bottomfraction}{.85}
\renewcommand{\textfraction}{.07}
\renewcommand{\floatpagefraction}{.65}
\captionsetup{font=small,labelfont=bf}
\newcommand{\kms}{\mathrm{km\,s^{-1}}}
\newcommand{\Msol}{\mathrm{M}_{\odot}}
\newcommand{\E}{\mathbb{E}}
\newcommand{\ind}{\mathbf{1}}
\newcommand{\figfile}[2]{\IfFileExists{figures/#1.pdf}{\includegraphics[width=#2\linewidth]{figures/#1.pdf}}{\fbox{\parbox[c][35mm][c]{.88\linewidth}{\centering Figure #1}}}}
\newcommand{\tallfigfile}[1]{\IfFileExists{figures/#1.pdf}{\includegraphics[width=\linewidth,height=.72\textheight,keepaspectratio]{figures/#1.pdf}}{\figfile{#1}{.94}}}
% Edit authorship here before submission. Do not infer affiliations from a repository.
\newcommand{\PaperAuthors}{Nalin Dhiman}
\newcommand{\PaperAffiliation}{}


\title{Measuring spiral structure and testing its incremental information for galaxy redshift estimation}
\author{\PaperAuthors\PaperAffiliation}
\date{}
\begin{document}
\maketitle
\begin{abstract}
Resolved spiral structure is an accessible source of information about galaxies,
but its usefulness for locating galaxies must be measured separately from the
ability to recognize their arms. We present a compact, interpretable pipeline
that combines calibrated three-band images, multiscale spatial filters and
human arm-region annotations, then tests whether the resulting structure
summaries improve catalog-redshift prediction. Sky groups separate training,
calibration and testing. On 29 test galaxies, the arm-region classifier reaches
mean average precision 0.720, compared with 0.655 for a capacity-matched
brightness/radius comparator; the equal-group gain is 0.0848, with a 95\%
bootstrap interval of 0.0475--0.1247. Applied without retraining to additional
strong spirals in previously excluded sky regions, it reaches 0.689 on 57
evaluable galaxies. Eight predicted-region summaries, however, do not establish
an incremental redshift gain over photometry, apparent geometry and simple image
descriptors. The equal-group reduction in absolute redshift error is
$-0.000471$ on the original 29 galaxies and $+0.000440$ on 59 additional galaxies;
both intervals include zero. Stronger retrospective morphology baselines leave
this conclusion unchanged. A separate human-centerline experiment shows why
region agreement cannot be used as a substitute for validated pitch or width
measurements. The pipeline supplies useful region measurements and a controlled
test of their downstream information, while keeping angular structure,
statistical redshift and physically calibrated distance distinct.
\end{abstract}
\noindent\textbf{Keywords:} galaxies: spiral --- methods: image processing ---
methods: statistical --- techniques: photometric --- distance scale

\section{Introduction}
Spiral arms are among the most recognizable structures in astronomical images.
Their winding, separation, contrast and colors can be measured without first
assigning a mass model to the galaxy. A practical measurement system could
therefore support studies of disk structure, provide regions for spatially
resolved spectroscopy, and offer additional predictors for statistical redshift
inference. These applications ask different questions. Recognizing where people
draw an arm is a region-classification problem; determining its centerline or
pitch is a geometric problem; estimating redshift is a supervised inference
problem. Reliable progress requires each link to be evaluated on its own.

Automatic arm measurement already has a substantial literature.
\citet{davis2014} introduced SpArcFiRe, which extracts arm segments and fits
logarithmic spirals. \citet{hart2017} combined these measurements with Galaxy Zoo
classifications to investigate arm formation. Citizen-science projects supply
complementary information: Galaxy Zoo 2 describes detailed visual morphology
\citep{willett2013}; Galaxy Zoo:3D supplies spatial masks of arms, bars and
foreground stars \citep{masters2021}; and Galaxy Builder provides multicomponent
models and human-guided arm traces \citep{lingard2020,lingard2021}. More recently,
the ZooBot:3D preprint describes large-scale neural arm/bar segmentation in
Legacy Survey images \citep{spindler2026}. These studies motivate a focused
question about the \emph{information contributed by an explicit arm representation},
rather than a claim to have invented arm detection.

Morphology can also help photometric-redshift estimation. The benefit depends
on the photometric bands, sample and representation. \citet{soo2018} studied
size, ellipticity, S\'ersic index and surface brightness, finding greater
improvements when fewer bands were available. \citet{pasquet2019} demonstrated
redshift prediction directly from multiband images. An arm-specific feature
may be useful even when global morphology is already supplied, but this must
be tested against a baseline with comparable model capacity and measurement
quality. Otherwise, an apparent improvement could reflect brightness, apparent
size, resolution or observational selection rather than spiral structure.

This paper makes three connected contributions. First, it presents a transparent
three-band measurement pipeline with angular filter scales, explicit masks and
separate artifacts for visualization and downstream inference. Second, it
quantifies region-ranking performance against human annotations, including a
spatially distinct additional cohort. Third, it measures the incremental value
of eight predicted-region summaries for redshift, with common regressors,
shared quality covariates and uncertainty calculated at the sky-group level.
Separate geometric and dynamical analyses establish which further physical
interpretations would require additional evidence. The software is intended
to make these distinctions operational, so that a visually convincing map
does not silently become a claimed physical measurement.

\section{What an image can measure}
\label{sec:theory}
\subsection{Angular structure, physical size and redshift}
For a small angular separation $\theta$, a transverse physical separation is
\begin{equation}
 L=D_A\theta, \qquad \theta\ \hbox{in radians},
 \label{eq:angular}
\end{equation}
where $D_A$ is angular-diameter distance \citep{hogg1999}. A digital separation
$s$ pixels becomes $\theta=sp$ for a local angular scale $p$. This determines an
angle, not a distance: the same angular pattern can arise from different
intrinsic sizes at different distances. Similarly, an apparent flux constrains
luminosity only through a distance-dependent conversion. In the nearby Euclidean
limit, simultaneous changes $D\mapsto\lambda D$, $L_{\rm size}\mapsto\lambda
L_{\rm size}$ and $L_{\rm luminosity}\mapsto\lambda^2L_{\rm luminosity}$ preserve
apparent size and total flux. Population relations, spectral information and
selection can make some such combinations more probable than others. A learned
redshift estimator uses those calibrated statistical relations.

Redshift and distance are consequently separate outputs. The target in this
study is a catalog spectroscopic redshift, $z$. A cosmological conversion would
require a cosmological model, and local peculiar velocities also affect the
relation between observed redshift and expansion distance. We report errors in
$z$ and do not reinterpret them as independently validated distance errors.
Physical arm separations can be reported when an external distance is supplied,
with that distance uncertainty propagated.

\subsection{Projected and deprojected spiral geometry}
Let $(x,y)$ be angular coordinates relative to a supplied center. Rotating by the
disk position angle gives $(x',y')$. For a thin circular disk inclined by $i$,
\begin{equation}
 X=x',\qquad Y=\frac{y'}{\cos i},\qquad
 R_\theta=(X^2+Y^2)^{1/2},\qquad \phi=\operatorname{atan2}(Y,X).
 \label{eq:deproject}
\end{equation}
The approximation becomes sensitive to geometry errors at high inclination.
It is appropriate only after specifying what defines the disk plane; a bar,
bulge or foreground object can bias a moment-based estimate.

A logarithmic arm segment obeys
\begin{equation}
 R_\theta(\phi)=R_0\exp[b(\phi-\phi_0)],\qquad
 |\psi|=\arctan |b|,
 \label{eq:pitch}
\end{equation}
where $\psi$ is the pitch angle relative to the local circular direction.
Pitch is dimensionless and does not require a distance, but its measured value
depends on the center, inclination, position angle and selected arm segment.
Different arms in one galaxy need not share one pitch
\citep{davis2014,lingard2021}. A useful region map alone does not determine these
geometric quantities. For this reason the primary classifier works in a fixed
\emph{projected circular annulus}; deprojection is reserved for a separate
centerline validation with externally supplied disk geometry.

\subsection{Passbands and spatial filters}
A passband measures a different part of the spectral energy distribution.
A spatial filter transforms measurements already present in an image. We use
the calibrated $g,r,z$ science planes individually, preserving their different
brightness and color information. A single-band image is naturally a grayscale
measurement. Converting a rendered color image into black and white changes its
display representation; it does not supply a new passband or restore calibrated
fluxes that were lost during rendering.

For a band image $I_b$, write its radial-profile residual as
$J_b=I_b-B_b-m_b(R_\theta)$, where $B_b$ is a robust background and $m_b$ a
radial median profile. Gaussian scale space is
\begin{equation}
 L_b(\bm{x};\sigma)=G_\sigma*J_b,
 \qquad
 G_\sigma(\bm{x})=\frac{1}{2\pi\sigma^2}
 \exp\left[-\frac{|\bm{x}|^2}{2\sigma^2}\right].
 \label{eq:scale}
\end{equation}
Differences of smoothed images emphasize local contrast; gradients locate
changes; Hessian eigenvalues distinguish elongated bright ridges from locally
isotropic structure. Scale-normalized derivatives follow standard scale-space
reasoning \citep{lindeberg1998}, and the ridge descriptor is related to established
Hessian-based enhancement \citep{frangi1998}. Our scales are fixed angular
standard deviations of 0.4, 0.8 and 1.6 arcsec. They are measurement choices,
not inferred physical arm widths. The full feature definitions and treatment
of noise, boundaries and masks appear in the Supplementary Material.

\section{Observations and experimental design}
\label{sec:data}
The original sample contains 160 resolved strong spirals selected
deterministically from the public Galaxy Zoo:3D catalog, excluding 32 earlier
feasibility objects. The spiral-vote fraction is at least 0.7 and the source
selection requests one center. Native Legacy Survey DR9 $g,r,z$ images,
inverse-variance planes and local point-spread functions (PSFs) provide the
science measurements \citep{dey2019}. Human spatial labels and centers come from
Galaxy Zoo:3D \citep{masters2021}. Observed catalog $ugriz$ photometry, apparent
geometry and the heliocentric \texttt{NSA\_Z} target are matched by identity to
MaNGA DRPall, within the SDSS data framework
\citep{york2000,bundy2015,dr172022}.

Ten-degree right-ascension/declination tiles define sky groups. There are
91 training galaxies in 38 groups, 40 calibration galaxies in 13 groups and
29 test galaxies in 13 groups. No group crosses these boundaries. Spatial
grouping addresses a common weakness of random splitting in structured data
\citep{roberts2017}; it is not a proof that all remaining groups have independent
calibration errors. Training is used for fitted coefficients and model
selection, calibration for a visualization sigmoid and prediction intervals,
and testing for the final original comparison.

An additional acquisition excludes all 192 earlier identifiers and all 67
earlier sky tiles. It contains 99 objects in 39 new tiles: 59 eligible strong
spirals in 30 groups, plus 40 weak-spiral or low-vote comparison objects. The
59 strong spirals form the principal additional redshift cohort. All are
evaluated with the original frozen models, without threshold or calibration
updates. This tests new sky regions within the same imaging survey.

Coverage is retained explicitly. One additional strong spiral has no usable
native science/PSF product after repeated requests. Another lacks both scoring
classes after masking. Consequently arm AP is defined for 57 of 59 strong
spirals; catalog-redshift predictions retain all 59 rows using training-fitted
imputation when needed. A low-vote or unasked arm task is treated as unknown,
not as a verified arm-free control.

\begin{figure}[htbp]
\centering\figfile{pipeline_splits}{.98}
\caption{Measurement and evaluation design. The image-to-region model and
the redshift regressors are trained using only training sky groups. Display
calibration is separate from the raw-score representation used for redshift.
The additional cohort is evaluated with the frozen transformations and models.}
\label{fig:pipeline}
\end{figure}

\section{Measurement and inference methods}
\label{sec:methods}
\subsection{Registered science measurements and reference regions}
Each galaxy is measured within a fixed 3--23 arcsec annulus about its supplied
human center. Foreground-star labels exclude contaminated pixels. All bands use
the same native grid: only the reference labels are mapped through WCS by
nearest neighbor. The reader checks the target-containing image brick, band,
image/inverse-variance pairing and coordinate registration. PSF matching broadens
each band to the worst local Gaussian moment-equivalent width. A shared guard
excludes pixels within four times the largest filter scale of masked holes or
boundaries, in addition to the convolution mask treatment.

The reference-positive class is a pixel covered by at least three human arm
polygons. Other scored pixels form the comparison class. This is an annotation
agreement target: an unmarked pixel can contain an incompletely traced arm.
Five-vote labels provide a sensitivity analysis, with the changing positive
prevalence and evaluable sample reported separately.

\subsection{A compact multiscale region classifier}
The 44 local features comprise three radius terms, 13 intensity/contrast/shape
terms in each of three bands and two signed color terms. They include radial
profile and residual brightness, differences of Gaussians, gradients, a
Laplacian-like response, bright-ridge strength and ridge orientation relative
to the center. None requires a redshift, physical size or mass model.

Seven fixed logistic/ExtraTrees candidates are compared by three-fold
sky-group cross-validation within the training set. ExtraTrees uses randomized
tree ensembles \citep{geurts2006}, implemented with scikit-learn
\citep{pedregosa2011}. The selected full model has 80 trees, maximum depth 12,
minimum leaf size 32 and a feature fraction of 0.7. Each training galaxy
contributes at most 400 positive and 400 comparison pixels, with equal total
weight per galaxy. Thus a galaxy with more image pixels cannot dominate simply
because it covers more area.

The brightness/radius comparator has the same eligible logistic/tree capacities
but only three radius terms, three normalized band intensities and the $g$-band
radial profile. Training cross-validation also selects its depth-12 tree model.
Raw $g$ intensity and a circular-profile score provide simpler comparisons.
This design asks whether local multiscale features add beyond intensity and
radius under a controlled learner budget.

The visualization artifact applies a positive-slope sigmoid to the logit of
the tree score, following the calibration approach associated with
\citet{platt1999}. Calibration data select one display threshold, 0.30.
Average precision (AP) evaluates score ranking and is unaffected by the monotone
calibration. We compute AP separately for each galaxy and average it, avoiding
a pooled-pixel metric dominated by large objects. Fixed-threshold precision,
recall, F1 and intersection-over-union are secondary region diagnostics.

\subsection{A separate representation for downstream redshift}
For downstream inference we save the training-only classifier \emph{without}
the display sigmoid, using a fixed raw threshold of 0.5. Training galaxies
receive descriptors from a classifier fitted without their entire sky group;
calibration and test galaxies receive descriptors from the complete training-only
model. The summary annulus depends on the image masks, not on human arm coverage.
Calibration arm labels therefore do not enter redshift features.

Let $s_j\in[0,1]$ be this raw score at valid pixel $j$ and $u_j$ its radius
divided by 23 arcsec. The eight arm summaries are the mean, 90th percentile and
standard deviation of $s$, the score-weighted mean of $u$, mean scores inside
and outside $u=0.55$, mean binary entropy and the fraction with $s\geq0.5$.
The raw score ranks agreement with the annotation protocol; balancing classes
during training means it is not a calibrated physical arm probability.

Four feature families share one regression design. P contains five observed
magnitudes and four adjacent colors. PG adds angular half-light radius and axis
ratio. PGI adds seven simple image descriptors: three annular positive fluxes,
a radial half-flux scale, residual contrast and two colors. PGIA adds the eight
arm summaries. All share magnitude-error, PSF-width and image-noise covariates,
with consistent missing-feature handling. This reduces the opportunity for arm
availability or resolution to act as an exclusive proxy for redshift.

Each regressor uses 128 ExtraTrees, minimum leaf size four and the same seed,
with training-median imputation and missingness indicators. The fitted target is
$t=\log(1+z)$ and predictions are transformed back through
$\hat z=\exp(\hat t)-1$. Training cross-validation selects P as the retained
deployment family. The scientific arm contrast remains PGI versus PGIA,
regardless of this choice. No original or additional test result changes the
deployment selection.

\subsection{Paired effects and uncertainty}
For a sky group $g$ containing $n_g$ galaxies, define
\begin{equation}
 d_g=\frac{1}{n_g}\sum_{i\in g}
 \left[|z_i-\hat z_{i,\rm PGI}|-|z_i-\hat z_{i,\rm PGIA}|\right],
 \qquad \Delta=\frac{1}{G}\sum_{g=1}^G d_g.
 \label{eq:effect}
\end{equation}
Positive $\Delta$ favors the arm summaries. Equal-group weighting differs
deliberately from galaxy-weighted mean absolute error (MAE). Percentile
bootstrap intervals resample whole groups \citep{efron1979}; the uncertainty
unit is not a pixel. The two cohorts are reported separately. The final
statistical analysis uses 100,000 group resamples and also examines a
Student-$t$ interval as a small-sample sensitivity.

For interpretable effect sizes we additionally examine a reduction
$\delta_z=0.001$. This is a retrospective practical margin for this pilot,
not a physically validated universal threshold. A confidence interval that
crosses zero is inconclusive about the sign. To support practical equivalence
at this margin, the relevant interval would need to be contained within
$[-\delta_z,\delta_z]$; a nonsignificant difference alone is insufficient.

Nominal 90\% split-conformal intervals use the maximum absolute residual in
$t$ within each calibration sky group \citep{lei2018,angelopoulos2021}.
With 13 groups the finite-sample rank is the largest calibration group residual.
These intervals concern catalog redshift in comparable exchangeable groups.
Distribution shift requires separate treatment \citep{tibshirani2019}.

\section{Results}
\label{sec:results}
\subsection{Multiscale features improve agreement with arm regions}
The original test classifier reaches galaxy-macro AP 0.7197, versus 0.6551
for the capacity-matched brightness/radius comparator, 0.6267 for raw $g$
brightness and 0.4904 for the circular-profile score. The improvement is not
limited to a visibly selected example: the equal-sky-group AP gain over the
matched comparator is 0.0848, with a 95\% bootstrap interval of
0.0475--0.1247. At the fixed display threshold, macro precision is 0.590,
recall 0.678 and F1 0.581. These results establish a useful ranking improvement
for the specified region task and selected population.

Without retraining, mean AP is 0.6890 on the 57 evaluable additional strong
spirals, with a sky-group bootstrap interval of 0.6409--0.7376. Median
fixed-threshold F1 is 0.579. This provides evidence that the region-ranking
performance extends to additional sky regions in the same survey. It does
not validate every output as a complete physical arm.

\begin{figure}[p]
\centering\tallfigfile{arm_region_examples}
\caption{Fixed science-image, raw-score and reference panels: the first two
additional strong spirals, first weak/bar-proxy object and first unasked control
in the pre-scoring gallery. Scores use the training-only uncalibrated branch;
they are distinct from geometric centerlines or physical arm probabilities.
The no-consensus case is retained, and an unasked task has unknown truth.
White valid unmarked pixels are comparison labels, not verified physical absence.
Legacy Surveys / D. Lang (Perimeter Institute). Human region references:
SDSS-IV and Galaxy Zoo:3D volunteers \citep{masters2021}.}
\label{fig:examples}
\end{figure}
\begin{figure}[htbp]
\centering\figfile{arm_region_validation}{.94}
\caption{Arm-region validation. Original-cohort comparisons use the same
scoring pixels and independently selected matched-capacity baseline; additional
cohort performance uses the frozen full classifier. Galaxy-macro AP and
equal-sky-group contrasts are distinct estimands. Reference-positive pixels
require at least three human polygons.}
\label{fig:ap}
\end{figure}

\FloatBarrier
\subsection{The incremental redshift effect remains unresolved}
Table~\ref{tab:redshift} reports galaxy-weighted MAE for all four frozen
families. The original-cohort PGI-to-PGIA contrast is negative at equal-group
weighting, $\Delta=-0.000471$, with a 95\% interval
$[-0.001391,0.000367]$. On the additional 59 strong spirals it is positive,
$\Delta=0.000440$, with interval $[-0.000555,0.001452]$.
Both include zero. The original interval excludes a gain as large as the
retrospective margin of 0.001, while the additional-cohort interval still
permits a gain at that margin. Neither establishes practical equivalence.

\begin{table}[tb]
\centering
\caption{Catalog-redshift prediction by the original frozen regressors.
MAE weights galaxies equally; the principal arm contrast in
equation~\eqref{eq:effect} weights sky groups equally. All 59 additional rows
remain in the analysis, including the missing-image row.}
\label{tab:redshift}
\begin{tabular}{lrr}
\toprule
Feature family & Original: 29 galaxies & Additional: 59 galaxies\\
\midrule
P: photometry & 0.009635 & 0.013190\\
PG: P + apparent geometry & 0.009469 & 0.012560\\
PGI: PG + simple image summaries & 0.009266 & 0.012767\\
PGIA: PGI + arm-region summaries & 0.009751 & 0.012532\\
\bottomrule
\end{tabular}
\end{table}

\begin{figure}[htbp]
\centering\figfile{grouped_redshift_effects}{.95}
\caption{Paired reductions in absolute catalog-redshift error from adding
arm-region summaries. Positive values favor arms. The original and additional
cohorts use frozen original models. L5/L3 comparisons are retrospective
controls. Intervals resample whole sky groups. The shaded practical margin
$\pm0.001$ was introduced retrospectively and does not convert this pilot into
a confirmatory equivalence study.}
\label{fig:effects}
\end{figure}

The P family, retained by training selection, has original coverage 29/29
with median redshift-interval width 0.08954. On the additional cohort, coverage
is 57/59 and median width 0.09043. The corresponding whole-group coverage is
28/30. These broad intervals are informative about limited precision, even
when empirical coverage is high. Input-only support checks flag 15/59
additional P predictions for possible abstention. Table~\ref{tab:redshift}
keeps all rows instead of reporting only favorable accepted predictions.

\begin{figure}[htbp]
\centering\figfile{redshift_predictions}{.94}
\caption{Frozen PGI and PGIA predictions and absolute-error cumulative
distributions for all 59 additional strong spirals. The missing-image row is
imputed and marked, and no input-support screening removes rows. The
galaxy-weighted MAEs differ from the equal-group contrast in
Figure~\ref{fig:effects}. These are catalog redshifts, not independent geometric
distances.}
\label{fig:predictions}
\end{figure}

\FloatBarrier
\subsection{Stronger photometric morphology controls}
We performed an additional, explicitly retrospective baseline analysis
motivated by \citet{soo2018}. L5 contains observed five-band photometry,
apparent geometry, observed $r$-band surface brightness and S\'ersic index,
plus the shared quality block. L5A adds the same eight arm summaries. L3/L3A
restrict explicit photometry and magnitude errors to $g,r,z$, retaining the
specified angular structure and $grz$ image-quality information. All use the
same 91 training galaxies and common ExtraTrees capacity; the additional
cohort does not select new hyperparameters.

On 59 additional strong spirals, equal-group MAE falls from 0.013132 to
0.012547 for L5/L5A, a gain of 0.000584 with interval
$[-0.000331,0.001658]$. For L3/L3A it falls from 0.012928 to 0.012786,
a gain of 0.000142 with interval $[-0.000802,0.001195]$. These controls leave
the incremental arm result unresolved. They are a controlled implementation
of literature-motivated observable features, not a reproduction of the
Morpho-z algorithm, training population or survey. The three-band restriction
also does not constitute a cross-survey experiment.

\subsection{Centerlines and pitch require their own validation}
A separate 24-object Galaxy Builder sample supplies genuine human centerlines
and fitted disk geometry \citep{lingard2020,lingard2021}. Twelve development
objects select one classical logarithmic-template configuration; twelve test
objects assess localization within the supplied PSF tolerance. Test mean
centerline F1 is 0.223, with a whole-galaxy interval of 0.105--0.353. The
extractor predicts 70 segments against 19 retained reference segments; seven
spatially matched segments have median absolute pitch error 11.42 degrees.
This fails the development goal of F1 at least 0.7 and matched pitch error
at most 5 degrees. It is a different experiment from the supervised region
classifier and should not be interpreted as its centerline accuracy.

\begin{figure}[htbp]
\centering\figfile{centerline_validation}{.93}
\caption{Separate human-centerline validation of the classical geometry
prototype. Region labels and centerlines answer different questions. Pitch
error is conditional on the seven matched segments, while localization F1
includes unmatched predictions. SDSS imaging; Galaxy Builder volunteers and
author-fitted references \citep{lingard2020,lingard2021}.}
\label{fig:centerlines}
\end{figure}

An archival comparison uses actual published SpArcFiRe-related galaxy pitches
from the \citet{hart2017} catalog. All 24 Galaxy Builder identifiers match;
13 have finite reliable-arc pitches and 11 retain the published missing
sentinel. On the six measured test identities, the mean absolute difference
from the Galaxy Builder reference summary is 7.73 degrees, with a conditional
interval of 5.06--10.69 degrees. The two galaxy summaries weight different
selected arcs, and this is a catalog comparison rather than a rerun of the
author pipeline. It supports attention to reference conventions and missing
coverage, not a ranking of the present extractor against SpArcFiRe.

\FloatBarrier
\section{Interpretation and applications}
The clearest positive result is that a compact multiscale representation
improves region ranking beyond brightness and radius under a matched learner
budget. Its constituent measurements have a direct use: arm-associated pixels
can be compared with interarm regions, cross-matched to resolved spectroscopy,
or used to prioritize objects for geometric modeling. The frozen additional
cohort demonstrates useful same-survey transfer of that region task.

The redshift experiment addresses a further question, and its answer is more
limited. Once colors, apparent geometry and quality information are present,
the eight regional summaries have no established additional value in these
cohorts. This does not negate the region measurement. Nor does it bound all
information in an image: the summaries may discard topology, continuous
centerlines, segment-to-segment pitch variation or stellar-population structure.
The result concerns a defined representation and regression family.

The physical distinction is equally useful. A rotation curve constrains the
gravitational field through $v_c^2=R\,\partial_R\Phi$; an image measures
projected light. Combining them requires an explicit light-to-mass calibration,
disk geometry and independently measured velocities. Established Tully--Fisher
methods illustrate this additional information \citep{tully1977,schombert2020}.
In the Supplementary Material we provide a finite pressure-supported equilibrium
calculation, derived from known lowered-isothermal models
\citep{king1966,gieles2015}, to keep dynamical interpretations physically
consistent. This is a separate conditional model, not a source of missing
distance labels or a viscosity measurement from pixels.

\section{Limitations and a focused route forward}
The samples contain resolved, strongly selected nearby spirals. Human centers,
foreground masks and classifications remain external inputs. The additional
cohort changes sky regions within the same survey; it does not establish
transfer to another instrument, arbitrary morphology or unresolved galaxies.
Incomplete human regions define a reproducible annotation task, but not
complete physical truth or a verified real arm-free negative population.

Gaussian moment matching approximates the PSF and does not fully account for
non-Gaussian wings or noise covariance. Fixed angular filters probe different
physical scales across the sample. The finite annulus omits central and outer
structure, while star/border guards can remove substantial areas. These choices
are measured coverage conditions rather than proof of complete arm recovery.
Observed width accuracy has not been validated independently.

Statistical resolution is set by the 13 and 30 test groups, not the number
of pixels. Conditional bootstrap intervals omit retraining uncertainty,
annotation systematics and uncertainty in the selected population. The cached
out-of-fold representation protects training descriptors from their own group,
but the training-only choice of upstream representation is not fully nested
within every downstream cross-validation fold. Practical margins, the stronger
baselines and the archival comparison were introduced after initial results;
they should guide a new study rather than be described as prospective tests.

A next confirmatory study should fix its feature representation, survey
population, practical margin and missing-data policy before examining new
redshifts. It should include a genuine same-input comparison to mature arm
methods, with training-overlap checks for pretrained models such as ZooBot:3D
\citep{spindler2026}. A useful geometric extension must first pass independent
centerline and width tests. Conditional planning from the observed group
dispersion suggests roughly 66 independent groups for 80\% power to detect a
true 0.001 mean gain against zero, rising to 106 under an upper-variance
scenario. These calculations assume stable frozen models and comparable groups;
they are not evidence that such a gain exists or a promise of sample sufficiency.

\section{Conclusions}
The three-band multiscale classifier provides a measurable improvement in
human arm-region agreement over a matched brightness/radius baseline, and
retains useful performance on additional sky regions. This establishes a
practical region-analysis component for resolved spirals. Its eight regional
summaries do not yet establish an improvement in catalog-redshift prediction;
the original and additional grouped contrasts remain statistically inconclusive,
including after stronger observable morphology controls. Separate centerline
results show that geometric accuracy requires additional development and
independent validation. The resulting framework supports careful measurement
and calibrated inference while preserving the distinction between apparent
structure and physical distance.

\section*{Data and software availability}
The companion source distribution at
\url{https://github.com/nalin-dhiman/Galaxy-arm-project}
contains the publicly available measurement code, frozen-model
metadata, numerical result tables, figure-generation provenance and editable
\LaTeX\ sources. Original survey products are obtained from their credited
public archives; access and reuse follow the relevant archive terms. The
Supplementary Material describes the output schema and reproducibility sequence.

\section*{Acknowledgments}
This work uses the DESI Legacy Imaging Surveys, SDSS-IV/MaNGA, Galaxy Zoo:3D,
Galaxy Builder and SPARC, and acknowledges their scientific teams and volunteer
contributors. Image credit: Legacy Surveys / D. Lang (Perimeter Institute).
The Legacy Surveys and SDSS collaboration acknowledgment and citation pages
are \url{https://www.legacysurvey.org/acknowledgment/} and
\url{https://www.sdss4.org/collaboration/citing-sdss/}. Survey-specific credits
are retained with the derived figures.

\bibliographystyle{plainnat}
\bibliography{references}
\end{document}
